% Vietnamese translation for OI Public Library
% Source-File: IOI中国国家候选队论文/2005/潘震皓/潘震皓.ppt
\documentclass[11pt]{article}
\usepackage{oipl-vn}

\newcommand{\Oh}{\mathcal{O}}

\title{Bản trình chiếu: Lũy thừa nhanh trên nhóm hoán vị}
\author{Pan Zhenhao}
\date{2005}

\begin{document}
\maketitle

\origtitle{Fast exponentiation in permutation groups}
\sourcepath{IOI China candidate team papers / 2005 / Pan Zhenhao / Pan Zhenhao.ppt}

\section{Phạm vi bản dịch}

Bản PowerPoint đi cùng bài viết về lũy thừa trong nhóm hoán vị. Text stream
giữ lại các mốc CEOI 1998, ví dụ chu trình, các trường hợp
\(\bmod 3=0,1,2\), điều kiện \(\gcd(n,k)>1\), nhắc tới ma trận hoán vị và so
sánh
\[
  \Oh(n^3\log k)\Rightarrow \Oh(n+k),\qquad \Oh(nm+km).
\]
Bản ghi chú này đi theo bài viết đã dịch, nhưng nén lại thành mạch slide:
phân tích chu trình, tính lũy thừa nguyên, rồi nhìn sang bài toán khai căn
hoán vị.

\section{Hoán vị và chu trình}

Một hoán vị có thể phân tích thành các chu trình rời nhau. Nếu
\[
  T=(a_0\ a_1\ \ldots\ a_{\ell-1}),
\]
thì áp dụng \(T\) một lần đưa \(a_i\) sang \(a_{i+1}\). Áp dụng \(T^k\) nghĩa
là đi \(k\) bước trong chu trình:
\[
  a_i\mapsto a_{i+k\bmod \ell}.
\]
Vì vậy cấu trúc của \(T^k\) phụ thuộc vào \(\gcd(\ell,k)\). Nếu
\(\gcd(\ell,k)=1\), chu trình vẫn là một chu trình; nếu
\(\gcd(\ell,k)=g\), chu trình tách thành \(g\) chu trình con, mỗi chu trình đi
qua các vị trí cùng lớp dư modulo \(g\).

\section{Tính lũy thừa nguyên}

Cách trực tiếp có thể mô hình hóa hoán vị bằng ma trận hoán vị và dùng lũy
thừa nhanh, dẫn đến \(\Oh(n^3\log k)\). Nhưng với biểu diễn chu trình, mỗi
phần tử chỉ cần được thăm một lần. Với từng chu trình, ta đánh dấu các phần
tử chưa xử lý, bắt đầu từ một phần tử, rồi liên tục nhảy \(k\) vị trí cho đến
khi quay lại. Các phần tử thu được tạo thành một chu trình trong hoán vị
đích.

Ví dụ, nếu một chu trình có độ dài \(6\), thì:
\[
  \gcd(6,2)=2,\quad \gcd(6,3)=3,\quad \gcd(6,5)=1.
\]
Tương ứng, \(T^2\) tách thành hai chu trình, \(T^3\) tách thành ba chu trình,
còn \(T^5\) vẫn là một chu trình. Các slide dùng các trường hợp
\(\bmod 3=0,1,2\) để minh họa cách lớp dư quyết định nhóm phần tử.

Khi lưu tất cả chu trình bằng hai mảng, một mảng chứa dữ liệu liên tiếp và một
mảng chứa vị trí bắt đầu/độ dài, lũy thừa nguyên được tính trong thời gian
\(\Oh(n)\). Nếu cần kết hợp thêm một chu kỳ ngoài dài \(k\), ta có dạng
\(\Oh(n+k)\) như slide ghi.

\section{Ví dụ CEOI 1998}

Mốc CEOI 1998 trong slide liên quan đến bài có nhiều hoán vị được áp dụng theo
chu kỳ. Nếu mỗi năm chọn một hoán vị trong một danh sách tuần hoàn, trạng thái
sau \(y=xp+r\) năm có thể viết bằng một hoán vị \(A^x\) ghép với phần đầu
chu kỳ \(C_r\). Bài toán trở thành: tìm \(x\) nhỏ nhất sao cho một phương
trình hoán vị đúng.

Khi phân rã \(A\) và \(C_r\) theo chu trình, mỗi chu trình cho một điều kiện
đồng dư kiểu
\[
  x\equiv x_i\pmod {n_i}.
\]
Nếu có nhiều chu trình, ta cần kiểm tra tính tương thích của hệ đồng dư. Đây
là chỗ nhóm hoán vị gặp số học: thuật toán không còn là nhân hoán vị lặp đi
lặp lại, mà là giải các ràng buộc modulo sinh ra từ chu trình.

\section{Khai căn hoán vị}

Phần sau của bài viết nói về lũy thừa phân số, tức tìm \(S\) sao cho
\[
  S^k=T.
\]
Nếu độ dài chu trình nguyên tố cùng nhau với \(k\), ta đảo ngược cách nhảy
\(k\) bước để dựng \(S\). Nếu không nguyên tố cùng nhau, một chu trình trong
\(S\) khi nâng lũy thừa sẽ tách thành nhiều chu trình trong \(T\). Vì vậy để
khai căn, ta có thể phải ghép nhiều chu trình cùng độ dài của \(T\) thành một
chu trình lớn rồi mới dựng nghiệm.

Điều kiện cốt lõi là phải có đủ số chu trình cùng độ dài để ghép, thường gắn
với \(m=\gcd(\ell,k)\). Nếu thiếu chu trình, không tồn tại nghiệm. Nếu đủ, ta
ghép xen kẽ các chu trình rồi dùng lại thuật toán nhảy vị trí.

\section{Kết luận}

Slide nhấn mạnh rằng nhóm hoán vị giúp tránh mô hình ma trận nặng nề. Khi nhìn
đúng bằng chu trình, lũy thừa nguyên là tuyến tính, còn khai căn hoán vị trở
thành bài toán ghép/tách chu trình có điều kiện \(\gcd\). Đây là ví dụ điển
hình của việc thay biểu diễn dữ liệu để làm lộ cấu trúc toán học.

\end{document}
